%% ═══════════════════════════════════════════════════════════════════════════
\section{Lecture 4: Dimensional Analysis — From Quantum Atoms to Radiation and Flow}
%% ═══════════════════════════════════════════════════════════════════════════
\date{30/04/2026}

This lecture deepens our practice of dimensional analysis. We move from simple
mechanical examples through quantum mechanics (Bohr radius, positronium) to
electromagnetism (dipole radiation, Larmor formula) and fluid dynamics (drag
force, Reynolds number). A central theme is the power of \emph{identifying
which physical parameters are relevant} before building a dimensionless
combination.

%%───────────────────────────────────────────────────────────────────────────
\subsection{Review: Simple Dimensional Analysis Examples}
%%───────────────────────────────────────────────────────────────────────────

\textbf{The Physical Picture:}
In the previous lecture we introduced dimensional analysis. Here we warm up by
revisiting three examples that share the same structure: a single length scale,
a single acceleration, no mass — so there is a unique dimensionless combination.

\subsubsection*{The Simple Pendulum}
The relevant parameters are length $L$ and gravitational acceleration $g$.
Their dimensions are $[L] = \si{m}$ and $[g] = \si{m\,s^{-2}}$.  We want a
time $T$.

The only combination that yields seconds is:
\[
    T \propto \sqrt{\frac{L}{g}},
\]
so the period is $T = C\sqrt{L/g}$ for some dimensionless constant $C$ (which
turns out to be $2\pi$ from solving the equation of motion).

\subsubsection*{Free Fall from Height $H$}
A body dropped from height $H$ in gravity $g$. The relevant parameters are
exactly $H$ and $g$, so:
\[
    t_\text{fall} \sim \sqrt{\frac{H}{g}}.
\]
This is structurally identical to the pendulum — same dimensional reasoning, different physics.

\subsubsection*{Gravitational Free-Fall Timescale of a Star}
A star of mass $M$ and radius $r$ collapses under self-gravity.  Now we have
three parameters: $M$, $r$, and Newton's constant $G$.

\textbf{Dimensional check:}
\begin{align*}
    [G] &= \si{kg^{-1}\,m^3\,s^{-2}} \\
    [M] &= \si{kg}, \qquad [r] = \si{m}.
\end{align*}

To build a time we note that $[GM] = \si{m^3\,s^{-2}}$, so $GM/r^3$ has
dimensions $\si{s^{-2}}$, giving:
\[
    \boxed{t_\text{ff} \sim \sqrt{\frac{r^3}{GM}} = \frac{1}{\sqrt{G\bar{\rho}}}}
\]
where $\bar{\rho} = M/r^3$ is the mean density (up to a numerical factor of
order unity).

\noindent\textbf{Physical intuition:} The free-fall time depends only on the
mean density, not on the actual size or mass separately. A star with the same
density as water would fall on a timescale of hours; a white dwarf collapses in
seconds.

%%───────────────────────────────────────────────────────────────────────────
\subsection{The Bohr Radius from Dimensional Analysis}
%%───────────────────────────────────────────────────────────────────────────

\textbf{The Physical Picture:}
The hydrogen atom is a quantum-mechanical bound state of a proton and an
electron. The relevant physics is electromagnetic and quantum — there is no
classical length scale for the orbit. Dimensional analysis pinpoints the
single combination that sets the atomic size.

\textbf{Relevant parameters:}
\begin{itemize}
    \item Electron charge $e$ (in Gaussian units).  Instead of $e$ alone, it
    is convenient to use $e^2$ whose dimensions are
    $[e^2] = \si{erg\,cm} = \si{g\,cm^3\,s^{-2}}$.
    \item Electron mass $m_e$: $[m_e] = \si{g}$.
    \item Reduced Planck constant $\hbar$: $[\hbar] = \si{erg\,s} =
    \si{g\,cm^2\,s^{-1}}$.
\end{itemize}

\textbf{Why not the speed of light?}  The hydrogen ground state is
non-relativistic (the fine-structure constant $\alpha = e^2/\hbar c \approx
1/137 \ll 1$), so $c$ does not set the scale and we exclude it.

\textbf{Dimensional construction:}
We want $[a_0] = \si{cm}$.  We must eliminate grams and seconds.

\begin{enumerate}
    \item \emph{Eliminate seconds.}  $\hbar^2$ has dimensions
    $\si{g^2\,cm^4\,s^{-2}}$, and $e^2$ has dimensions
    $\si{g\,cm^3\,s^{-2}}$.  The ratio $\hbar^2/e^2$ has dimensions
    $\si{g\,cm}$ — one power of mass and one of length remain.

    \item \emph{Eliminate grams.}  Dividing by $m_e$ removes the gram:
    \[
        \frac{\hbar^2}{m_e\, e^2} \quad \Rightarrow \quad
        \left[\frac{\hbar^2}{m_e e^2}\right] = \si{cm}.
    \]
\end{enumerate}

This is the only combination with units of length, so:
\begin{equation}
    \boxed{a_0 = \frac{\hbar^2}{m_e e^2} \approx \SI{0.529}{\AA}}
    \label{eq:lec04:bohr_radius}
\end{equation}

\noindent\textbf{Physical intuition:} The Bohr radius is set by the competition between
the kinetic energy $\sim \hbar^2/(m_e r^2)$ (from the uncertainty principle,
$\Delta p \sim \hbar/r$) and the Coulomb potential energy $\sim e^2/r$.
Minimising their sum gives exactly $r = a_0$.

%%───────────────────────────────────────────────────────────────────────────
\subsection{Hydrogen-Like Atoms: Scaling with Nuclear Charge $Z$}
%%───────────────────────────────────────────────────────────────────────────

For an atom with nuclear charge $Ze$ and a single electron (hydrogen-like
ion), the only change to the dimensional analysis is $e^2 \to Ze^2$.  Therefore:
\begin{equation}
    a_0(Z) = \frac{\hbar^2}{m_e (Ze^2)/Z} = \frac{\hbar^2}{m_e e^2 Z}
    = \frac{a_0}{Z}.
    \label{eq:lec04:bohr_Z}
\end{equation}

The Bohr radius scales as $1/Z$.  The ground-state energy scales as:
\[
    E_Z = Z^2 E_H = Z^2 \times (\SI{-13.6}{eV}).
\]

\ex{Highly Ionised Iron: Fe$^{+25}$}{Iron has $Z = 26$; Fe$^{+25}$ is iron stripped of 25 electrons, leaving one
(the innermost).  Treating it as hydrogen-like:
\[
    a_0(\text{Fe}^{+25}) = \frac{a_0}{26} \approx \frac{\SI{0.529}{\AA}}{26}
    \approx \SI{0.020}{\AA}.
\]
The binding energy is:
\[
    E(\text{Fe}^{+25}) = 26^2 \times \SI{13.6}{eV} = 676 \times \SI{13.6}{eV}
    \approx \SI{9.2}{keV}.
\]
This is an X-ray transition, as expected for inner-shell electrons of heavy
elements.}

%%───────────────────────────────────────────────────────────────────────────
\subsection{Positronium: A Two-Body Bound State}
%%───────────────────────────────────────────────────────────────────────────

\textbf{The Physical Picture:}
Positronium (Ps) is a bound state of an electron ($e^-$) and a positron
($e^+$).  There is no heavy nucleus; both particles are equally light.  The
two-body problem reduces to a one-body problem with the \emph{reduced mass}.

\textbf{Reduced mass:}
\[
    \mu = \frac{m_e \cdot m_e}{m_e + m_e} = \frac{m_e}{2}.
\]

Replacing $m_e$ by $\mu$ in Eq.~\eqref{eq:lec04:bohr_radius}:
\begin{equation}
    a_0(\text{Ps}) = \frac{\hbar^2}{\mu\, e^2} = \frac{\hbar^2}{(m_e/2)\,e^2}
    = 2a_0 \approx \SI{1.06}{\AA}.
    \label{eq:lec04:bohr_Ps}
\end{equation}

Positronium is \emph{twice} as large as hydrogen.  The binding energy is
correspondingly halved:
\[
    E(\text{Ps}) = \frac{\mu}{m_e}\,E_H = \frac{1}{2}\,(\SI{-13.6}{eV})
    = \SI{-6.8}{eV}.
\]

\textbf{Decay of positronium:}
Positronium is unstable — the electron and positron annihilate.  The number of
photons in the final state is fixed by conservation laws.  Two cases arise:
\begin{itemize}
    \item \textbf{Para-positronium} (spins anti-parallel, $S=0$): decays to
    \textbf{2 photons}.  Energy conservation requires each photon to carry
    $m_e c^2 = \SI{511}{keV}$, giving two discrete lines in the spectrum.
    \item \textbf{Ortho-positronium} (spins parallel, $S=1$): charge-conjugation
    symmetry forbids 2-photon decay; decays to \textbf{3 photons} with a
    continuous spectrum summing to $2m_e c^2$.
\end{itemize}

\noindent\textbf{Physical intuition:} Para-positronium decays with a lifetime of
$\sim \SI{125}{ps}$; ortho-positronium lives $\sim \SI{142}{ns}$ — a factor
$\sim 1000$ longer because it must radiate three photons instead of two.

%%───────────────────────────────────────────────────────────────────────────
\subsection{Dipole Radiation: Deriving the Larmor Formula}
%%───────────────────────────────────────────────────────────────────────────

\textbf{The Physical Picture:}
An oscillating electric dipole radiates electromagnetic energy.  The
electric dipole moment is $p = qd$, where $q$ is the charge and $d$ the
separation.  When $p$ changes in time (e.g., an accelerating charge), energy
is radiated to infinity.  We want the radiated power $P$ from dimensional
analysis alone.

\textbf{Relevant parameters:}
We work in Gaussian units throughout.
\begin{itemize}
    \item Dipole moment $p$: $[p] = \si{esu\,cm} = \si{g^{1/2}\,cm^{3/2}\,s^{-1}}$.
    Equivalently, $[p^2] = \si{erg\,cm^3}$.
    \item The characteristic oscillation frequency $\omega$: $[\omega] = \si{s^{-1}}$.
    \item The speed of light $c$: $[c] = \si{cm\,s^{-1}}$.
\end{itemize}

\noindent\textbf{What parameters to exclude:} The charge $q$ and distance $d$
appear only in the combination $p = qd$. The medium is vacuum ($\epsilon_0 = 1$
in Gaussian units). No gravitational constant enters.

\textbf{Dimensional construction:}
We want $[P] = \si{erg\,s^{-1}} = \si{g\,cm^2\,s^{-3}}$.

\begin{enumerate}
    \item Start with $p^2$: dimensions $\si{g\,cm^3\,s^{-2}}$
    (energy $\times$ length).
    \item Multiply by $\omega^4$ to introduce extra $\si{s^{-4}}$:
    $[p^2 \omega^4] = \si{g\,cm^3\,s^{-6}}$.
    \item Divide by $c^3$: $[c^3] = \si{cm^3\,s^{-3}}$, so
    \[
        \left[\frac{p^2 \omega^4}{c^3}\right]
        = \frac{\si{g\,cm^3\,s^{-6}}}{\si{cm^3\,s^{-3}}} = \si{g\,s^{-3}}
        = \si{erg\,cm^{-2}\,s^{-1}}.
    \]
    That's an intensity (power per area), not power — one factor of $\si{cm^2}$ is missing.
    We need the total power, so we still need an area $\sim (c/\omega)^2 = \lambda^2/(4\pi^2)$
    — but that is a dimensionless-factor issue.  Including it:
    \[
        P \sim \frac{p^2 \omega^4}{c^3}.
    \]
\end{enumerate}

\noindent\textbf{Physical insight on the $\omega^4$ dependence:} Note that
$p\omega^2 \equiv \ddot{p}$ (the second time derivative of the dipole moment).
Therefore:
\begin{equation}
    \boxed{P \sim \frac{\ddot{p}^{\,2}}{c^3}}
    \label{eq:lec04:larmor}
\end{equation}
This is the \textbf{Larmor radiation formula} (up to a factor $2/3$ in SI units
or $2/(3c^3)$ in Gaussian).  Because $\ddot{p} = q\ddot{d} = q a$ for a
single charge, the Larmor formula states:
\[
    P = \frac{2e^2 a^2}{3c^3} \quad \text{(Gaussian)}.
\]

\textbf{Why is $P \propto \omega^4$?}  Each factor of $\omega$ comes from a
time derivative: one from $\dot{p}$, one from $\ddot{p}$.  Squaring gives
$\omega^4$.  This is the reason blue light scatters more than red in Rayleigh
scattering ($\lambda^{-4}$ dependence).

\textbf{Why must $P$ be even in $\omega$ (and in $q$)?}
By time-reversal symmetry (or charge-conjugation symmetry), the radiated power
cannot change sign when we reverse the oscillation direction or flip all
charges: hence $P \propto \ddot{p}^{\,2}$, not $\ddot{p}$.

%%───────────────────────────────────────────────────────────────────────────
\subsection{Monopole Radiation is Forbidden}
%%───────────────────────────────────────────────────────────────────────────

\textbf{The Physical Picture:}
Could a time-varying monopole moment $Q(t)$ radiate?  If it could, the
formula would be $P \sim \ddot{Q}^{\,2}/c$ (by dimensional analysis with one
fewer power of length).  However, \textbf{charge conservation} forbids this.

\textbf{Argument:}
Gauss's law in vacuum: the total charge $Q = \oint \mathbf{E}\cdot\mathrm{d}\mathbf{A}$
is fixed by the electric field on a distant surface.  If charge is conserved
in the volume, $\dot{Q} = 0$, so $\ddot{Q} = 0$ identically.  There is no
source for monopole radiation.

\textbf{Alternative viewpoint (multipole expansion):} A single moving charge
can always be decomposed, at large distances, as:
\[
    \text{field} = \underbrace{\text{monopole (static)}}_{\text{no radiation}} +
    \underbrace{\text{dipole (time-varying)}}_{\text{radiates}},
\]
because any motion of the charge creates a time-varying dipole moment
$p = q\,\mathbf{r}(t)$ while the monopole $q$ stays constant.  The leading
radiated term is therefore \emph{always} dipole, giving Eq.~\eqref{eq:lec04:larmor}.

%%───────────────────────────────────────────────────────────────────────────
\subsection{Gravitational Wave Radiation: Why It Is So Weak}
%%───────────────────────────────────────────────────────────────────────────

\textbf{Analogy with electromagnetism:}
In electromagnetism the hierarchy is:
\[
    \text{monopole (forbidden)} \to \text{dipole (leading)} \to \text{quadrupole (subleading)}.
\]
In gravity, by analogy, we might expect monopole and dipole radiation.  Both
are forbidden:
\begin{itemize}
    \item \textbf{Monopole forbidden:} mass-energy conservation means
    $\ddot{M} = 0$.
    \item \textbf{Dipole forbidden:} conservation of momentum means the
    dipole moment $\mathbf{D} = \sum m_i \mathbf{r}_i$ has $\ddot{\mathbf{D}}
    = \sum m_i \ddot{\mathbf{r}}_i = 0$ (Newton's third law / center-of-mass
    motion).
\end{itemize}

Therefore the \textbf{leading gravitational radiation is quadrupolar}:
\begin{equation}
    P_\text{GW} \sim \frac{G}{c^5} \left(\dddot{Q}_{ij}\right)^2
    \label{eq:lec04:GW}
\end{equation}
where $Q_{ij} = \sum m_a r_i^{(a)} r_j^{(a)}$ is the mass quadrupole tensor.
Two extra derivatives (relative to the dipole formula) mean two extra powers
of $\omega$, suppressing the radiation further.

\textbf{Why gravitational waves are so hard to detect:}
Compared to electromagnetic radiation:
\begin{enumerate}
    \item Radiation is suppressed by two extra powers of $\omega$ (quadrupole
    vs.\ dipole), so ordinary astrophysical sources radiate
    gravitational waves at a tiny rate.
    \item Gravitational wave detectors (LIGO/Virgo) couple to the wave via
    $G_N$, which is intrinsically tiny, making the transducer inefficient.
\end{enumerate}
Only extremely compact, rapidly-orbiting systems (binary black holes or
neutron stars) produce detectable waves.  The strain measured by LIGO is of
order $h \sim 10^{-21}$, meaning the \SI{4}{km} arms change length by
$\sim \SI{e-18}{m}$ — less than a proton radius.

%%───────────────────────────────────────────────────────────────────────────
\subsection{Drag Force and the Reynolds Number}
%%───────────────────────────────────────────────────────────────────────────

\textbf{The Physical Picture:}
A solid object of cross-sectional area $A$ moves at velocity $v$ through a
fluid of density $\rho$ and dynamic viscosity $\eta$.  The drag force $F_d$
depends on the regime of flow — turbulent or laminar — which is characterised
by a single dimensionless number.

\subsubsection*{Turbulent (High Reynolds Number) Limit}

In the turbulent limit, viscosity is irrelevant to the \emph{bulk} drag (even
though it controls the small-scale dissipation).  The relevant parameters are
$\rho$, $v$, and $A$:
\[
    [F_d] = \si{g\,cm\,s^{-2}},\quad [\rho] = \si{g\,cm^{-3}},\quad
    [v] = \si{cm\,s^{-1}},\quad [A] = \si{cm^2}.
\]
The unique combination is:
\begin{equation}
    \boxed{F_d = C_D\,\rho v^2 A}
    \label{eq:lec04:drag_turbulent}
\end{equation}
where $C_D$ is a dimensionless drag coefficient of order unity (depends on
shape).

\subsubsection*{The Reynolds Number}

To determine whether the flow is turbulent or laminar, we form a dimensionless
number from all the relevant parameters.  Including viscosity $\eta$ with
$[\eta] = \si{g\,cm^{-1}\,s^{-1}}$, the combination
\[
    \nu = \frac{\eta}{\rho}, \quad [\nu] = \si{cm^2\,s^{-1}},
\]
is the kinematic viscosity.  The single dimensionless parameter is:
\begin{equation}
    \boxed{\text{Re} = \frac{\sqrt{A}\,v}{\nu} = \frac{\rho\,\sqrt{A}\,v}{\eta}}
    \label{eq:lec04:reynolds}
\end{equation}
(using $\sqrt{A}$ as the characteristic length scale of the body).

\begin{itemize}
    \item $\text{Re} \gg 1$ \textbf{(turbulent)}: The inertia of the fluid
    dominates; energy cascades from large eddies to progressively smaller ones
    (Kolmogorov cascade) until it is dissipated at the viscous scale.  Drag is
    given by Eq.~\eqref{eq:lec04:drag_turbulent}.
    \item $\text{Re} \ll 1$ \textbf{(laminar / Stokes flow)}: Viscosity
    dominates; streamlines are smooth and symmetric.
\end{itemize}

\subsubsection*{Laminar (Stokes) Drag}

In the laminar limit, the drag force is linear in velocity (from the
Navier–Stokes equation linearised in $v$):
\[
    F_d \propto \eta\,v\,\sqrt{A} \propto \eta\,v\,R.
\]
For a sphere of radius $R$ this gives the exact Stokes result:
\begin{equation}
    F_d = 6\pi\,\eta\,v\,R.
    \label{eq:lec04:stokes}
\end{equation}

\textbf{How to see this from dimensional analysis:}
The general result in the laminar limit must take the form
\[
    F_d = C_D(\text{Re})\,\rho v^2 A,\quad \text{with } C_D \propto \text{Re}^{-1}
    \text{ as Re}\to 0.
\]
Substituting $C_D = K/\text{Re}$ and Re $= \rho v \sqrt{A}/\eta$:
\[
    F_d = \frac{K}{\rho v \sqrt{A}/\eta}\,\rho v^2 A = K\,\eta\,v\,\sqrt{A}.
\]
The factor $K$ is determined by the geometry; for a sphere $K = 6\pi$.

\noindent\textbf{Physical intuition:} In laminar flow, the energy is dissipated
by viscous shear in the boundary layer.  The momentum transferred per unit time
(= drag force) is proportional to the velocity gradient, which is $\sim v/R$,
times the area $R^2$ and the viscosity: $F \sim \eta (v/R) R^2 = \eta v R$.

%%───────────────────────────────────────────────────────────────────────────
\subsection*{Summary of Key Results}
%%───────────────────────────────────────────────────────────────────────────

\begin{center}
\renewcommand{\arraystretch}{1.4}
\begin{tabular}{lll}
\toprule
\textbf{System} & \textbf{Result} & \textbf{Key parameters} \\
\midrule
Stellar free-fall time & $t_\text{ff} \sim 1/\sqrt{G\bar\rho}$ & $G,M,r$ \\
Bohr radius & $a_0 = \hbar^2/(m_e e^2)$ & $\hbar, m_e, e$ \\
Positronium radius & $a_0(\text{Ps}) = 2a_0$ & reduced mass $\mu = m_e/2$ \\
Hydrogen-like radius & $a_0(Z) = a_0/Z$ & nuclear charge $Z$ \\
Larmor radiation & $P \sim \ddot{p}^{\,2}/c^3$ & $p, \omega, c$ \\
GW emission & $P_\text{GW} \sim G\dddot{Q}^2/c^5$ & quadrupole, suppressed \\
Turbulent drag & $F_d \sim \rho v^2 A$ & $\rho, v, A$ \\
Laminar drag (Stokes) & $F_d = 6\pi\eta v R$ & $\eta, v, R$ \\
Reynolds number & Re $= \rho v \sqrt{A}/\eta$ & determines flow regime \\
\bottomrule
\end{tabular}
\end{center}

%%───────────────────────────────────────────────────────────────────────────
\subsection*{Recommended Reading}
%%───────────────────────────────────────────────────────────────────────────

\begin{itemize}
    \item \textbf{Dimensional Analysis (general):} Barenblatt,
    \textit{Scaling, Self-Similarity, and Intermediate Asymptotics}, Ch.~1–3 —
    the definitive treatment; covers the Buckingham $\pi$ theorem rigorously
    and applies it to fluid mechanics and beyond.

    \item \textbf{Bohr radius and hydrogen atom:} Griffiths,
    \textit{Introduction to Quantum Mechanics} (3rd ed.), §4.2 — derives the
    energy levels and wavefunctions from first principles, making explicit the
    role of $a_0$.

    \item \textbf{Larmor formula and dipole radiation:} Griffiths,
    \textit{Introduction to Electrodynamics} (4th ed.), §11.2 — clear
    derivation of the Larmor formula and its generalisations.
    Jackson, \textit{Classical Electrodynamics} (3rd ed.), Ch.~9 — the
    authoritative reference for multipole radiation.

    \item \textbf{Reynolds number and fluid flow:} Batchelor,
    \textit{An Introduction to Fluid Dynamics}, §4.9 — derivation of the
    Reynolds number from the Navier–Stokes equation and discussion of the
    laminar-to-turbulent transition.

    \item \textbf{Gravitational waves:} Maggiore,
    \textit{Gravitational Waves: Theory and Experiments} (Oxford), Ch.~1 —
    excellent introduction to the quadrupole formula and why monopole/dipole
    radiation are absent.
    \href{https://www.ligo.org/science/Publication-GW150914/}{LIGO Scientific
    Collaboration, GW150914 (2016)} — the landmark detection paper.

    \item \textbf{Positronium:}
    \href{https://en.wikipedia.org/wiki/Positronium}{Wikipedia: Positronium}
    — good overview of the two states, their lifetimes, and decay modes.
\end{itemize}
